\section{Actual scalar modes and their gauge vertices}\label{sec:modes}

We now retain the light vector multiplets and calculate their interactions
with two massless bundle scalars. All fields in this section belong to the
actual ASD connection $\mathcal A=\mathcal A(a,b)$. Fix $0<\kappa_0<1$
and restrict to
\begin{equation}\label{eq:hierarchycone}
 0<b\le\kappa_0a,\qquad r^2=2(a^2+b^2).
\end{equation}
Constants may depend on $\kappa_0$ and the fixed compactification, but not
on $b/a$. Sobolev norms below concern spatial derivatives on the fixed K3;
they do not denote derivatives with respect to $a$ or $b$.

\subsection{Harmonic representatives of the external states}
In the harmonic coordinates of section~\ref{subsec:harmonic}, choose the
two Dolbeault-unit vectors
\begin{equation}\label{eq:plane}
 \beta_u:\ \delta X_3=\sigma_1/\sqrt2,\qquad
 \beta_v:\ \delta X_3=\sigma_2/\sqrt2,
\end{equation}
with all other components zero. The affine Hermitian $X_3$ plane satisfies
all three moment equations. Its tangents are perpendicular to the gauge
orbit and its quaternionic images at the Pauli point. The corresponding
real-unit one-forms at the reference connection are
\begin{equation}\label{eq:realstates}
 \mathsf u_0=(\beta_u-\beta_u^*)/\sqrt2,\qquad
 \mathsf v_0=(\beta_v-\beta_v^*)/\sqrt2.
\end{equation}
Transport their classes by $\Psi$, take actual ASD harmonic representatives,
and normalize the resulting ordered pair by the positive inverse square
root of its $2\times2$ Gram matrix. Denote these real-unit external states
by $\mathsf u,\mathsf v$. This fixes their frame as well as their plane.
Let $\Pi_1$ be reference ASD harmonic projection on real one-forms.

\begin{lemma}\label{lem:actualstates}
In the fixed-reference unitary Coulomb presentation, for every fixed
spatial Sobolev order $s$,
\begin{align}\label{eq:actualstates}
 \|\Pi_1\mathsf u-\mathsf u_0\|_{H^s}
 +\|\Pi_1\mathsf v-\mathsf v_0\|_{H^s}&\le C_s r^2,\\
 \|(1-\Pi_1)\mathsf u\|_{H^s}
 +\|(1-\Pi_1)\mathsf v\|_{H^s}&\le C_s r.\nonumber
\end{align}
These estimates are uniform on \eqref{eq:hierarchycone}.
\end{lemma}
\begin{proof}
The fixed-complement reconstruction of
\cite[Lemmas~2.1--2.2, Proposition~2.3 and Lemma~3.1]{wuebben15}
has a holomorphic map $\mathfrak a(y)$, its derivative $J_y$, and a
complementary Hermitian correction $\sigma(y)$ satisfying
\begin{equation}\label{eq:reconstructionjets}
 \mathfrak a(y)=y+O(r^2),\quad J_y=I+O(r),\quad
 H_1J_y=I,\quad \bar\partial_0^*J_y=0,\quad \sigma(y)=O(r^2).
\end{equation}
Here $H_1$ is Dolbeault harmonic projection; the estimates hold in each
fixed spatial norm. Finite balancing replaces $x$ by $y=e^Sx$, where
$\|S\|\le Cr^2$, and hence $y-x=O(r^3)$. Conjugation acts only on
multiplicity matrices. Thus $y$ is still supported in slots $X_1,X_2$,
whereas $w=e^S\beta_u$ or $e^S\beta_v$ is supported in $X_3$. In
particular, $w=\beta_{u,v}+O(r^2)$ and $w$ is flat-orthogonal to the
entire complex orbit $O_y=\operatorname{im}C_y$ of the parallel gauge
algebra. The exact complex moment equation makes $w$ an integrable tangent.

Minimizing first over the fixed complementary complex gauge space produces
the positive Hermitian form $B_y$ on harmonic vectors, with
$\|B_y-I\|\le Cr^2$ \cite[Proposition~2.4]{wuebben15}. The remaining full
gauge minimum is the $B_y$-distance to $O_y$. If $o$ is the $B_y$-projection
of $w$ onto $O_y$, then
\begin{equation}\label{eq:orbitvectorbound}
 (1-Cr^2)|o|^2\le B_y(o,o)=B_y(w,o)
             =(B_y-I)(w,o),\qquad |o|\le Cr^2|w|.
\end{equation}
This estimates the orbit vector, without estimating a Lie algebra element
whose action gives that vector. No inverse small orbit eigenvalue enters.

For $w-o$, the complementary minimizer $\zeta$ satisfies
$\|\zeta\|_{H^{s+1}}\le C_s r|w-o|$ by the fixed-gap normal equation
and elliptic regularity. Its harmonic representative in the holomorphic
frame is
\[
 \gamma=J_y(w-o)-D_{\mathfrak a}\zeta,
 \qquad D_{\mathfrak a}=\bar\partial_0+[\mathfrak a,\cdot].
\]
It is closed by integrability and orthogonal to the full complex gauge
image for the HE metric, so it is Dolbeault harmonic for that metric.
Since $H_1J_y=I$ and $H_1\bar\partial_0\zeta=0$,
\[
 H_1\gamma=w-o-H_1[\mathfrak a,\zeta]
           =\beta_{u,v}+O(r^2).
\]
The complementary component is $O(r)$ in $H^s$.

Transport by $e^\sigma$ to the reference unitary bundle changes these
estimates only by $O(r^2)$. Differentiating this frame along a parameter,
or differentiating $S$, adds a complex gauge direction, removed by the
full harmonic minimum. No bound for either parameter derivative is used.
Passing to the real representatives \eqref{eq:realstates}, harmonic
orthogonality gives Gram matrix $I+O(r^2)$. Its inverse square root
preserves the estimates. Finally the fixed Coulomb correction is generated
by a complementary unitary gauge of size $O(r^2)$: the first connection
jet is already co-closed, and the reference complementary gauge Laplacian
has a fixed gap. Its conjugation changes these one-forms by $O(r^2)$.
\end{proof}

\subsection{The light eigenvector}
Put $D=d_{\mathcal A}$ and $\mathcal L=D^*D$. At the reference point let
\begin{equation}\label{eq:pauliprofiles}
 k_i=I_A\otimes i\sigma_i/\sqrt{6V_0},\qquad i=1,2,3,
\end{equation}
extended by zero on the other orthogonal-adjoint blocks. These are
orthonormal. The other 93 parallel sections form an exact fixed kernel
throughout the family with $a,b>0$.

\begin{lemma}\label{lem:lightdirection}
The normalized light eigenvector $p_L$ of $\mathcal L$ can be chosen with
$\|p_L-k_3\|_{H^s}\le C_s r^2$. The two heavier low eigenvectors span
a plane within $O(r^2)$ of $\operatorname{span}\{k_1,k_2\}$, uniformly
on \eqref{eq:hierarchycone}.
\end{lemma}
\begin{proof}
Write $\mathcal A=\mathcal A_0+\alpha$ in fixed Coulomb gauge. The
quadratic term in its unitary reconstruction is reference-harmonic
orthogonal: its $(0,1)$ part is $\mathfrak a_2-\bar\partial_0\sigma_2$,
with $\mathfrak a_2$ coexact. Finite balancing changes $x$ only by
$O(r^3)$. The complementary Coulomb correction has size $O(r^2)$;
its linear exact term has zero harmonic projection, and its commutator
with the first jet has size $O(r^3)$. Consequently
\begin{equation}\label{eq:connectiontwojet}
 \Pi_1\alpha=\alpha_1(x)+O(r^3),\qquad
 (1-\Pi_1)\alpha=O(r^2).
\end{equation}
These are value estimates in fixed spatial norms.

Let $\Pi_0$ project onto the full reference parallel algebra. Coulomb
orthogonality gives
$Q_0\mathcal L\Pi_0=Q_0[\alpha,\cdot]^*[\alpha,\cdot]\Pi_0=O(r^2)$.
The fixed complementary gap, or the eigenvector equation used in
\eqref{eq:lowcomplement}, represents the low cluster as a graph
$X:\operatorname{im}\Pi_0\to\operatorname{im}Q_0$ with
$X=O(r^2)$ in each fixed spatial norm. Its canonical isometry is
$U=(I+X)(I+X^*X)^{-1/2}$. Parallel commutators preserve reference
harmonic one-forms and their orthogonal complement. Equation
\eqref{eq:connectiontwojet} therefore gives
\begin{equation}\label{eq:effectivepauli}
 U^*\mathcal LU=\Lambda_x+O(r^4),\qquad
 \Lambda_x\big|_{\langle k_1,k_2,k_3\rangle}
 =\frac8{3V_0}\operatorname{diag}(a^2,a^2+b^2,b^2).
\end{equation}
Indeed the harmonic first-jet/quadratic-jet cross term vanishes, the
remaining diagonal error is $O(r^4)$, and the cross block times $X$
and the complementary graph energy have the same order.
The fixed 93-dimensional kernel decouples exactly. In its orthogonal
three-dimensional complement the light eigenvalue of $\Lambda_x$ has
separation at least $c_{\kappa_0}a^2$ from the other two. The error in
its unit eigenvector is bounded by $Cr^4/a^2\le Cr^2$. The orthogonal
plane and spatial estimates follow from the graph construction.
\end{proof}

The additive matrix estimate \eqref{eq:effectivepauli} does not prove a
relative bound for the light mass when $b\ll a^2$. That bound remains
the separate energy-relative conclusion of Proposition~\ref{prop:vectors}.

\subsection{Uniform cubic vertices}
Let $I_\alpha$, $\alpha=1,2,3$, be the parallel complex structures on
real one-forms. For an adjoint zero-form $p$ define
\begin{equation}\label{eq:vertices}
 j_p=\langle[p,\mathsf u],\mathsf v\rangle,\qquad
 b_p^\alpha(w,z)=\langle I_\alpha[p,w],z\rangle,
 \quad w,z\in\{\mathsf u,\mathsf v\}.
\end{equation}
The inner product uses $g_0$ and $\BF$. The second form is symmetric
because $I_\alpha$ and $\operatorname{ad}p$ are commuting skew-adjoint
operators. These are respectively the vector current and the Hessian of
a component of the gauge moment map, with signs fixed in
section~\ref{sec:scattering}.

\begin{proposition}\label{prop:vertices}
For the three low positive vector profiles,
\begin{equation}\label{eq:lowvertices}
 j_{p_L}=-\frac2{\sqrt{6V_0}}+O(r^2),\qquad
 j_{p_M},j_{p_H}=O(r^2),\qquad
 b_{p_j}^\alpha(w,z)=O(r^2).
\end{equation}
Every persistent zero mode has both vertices exactly zero. For a complete
orthonormal zero-form eigenbasis $\{p_n\}$,
\begin{equation}\label{eq:parsevalvertices}
 \sum_n\left[|j_{p_n}|^2+
 \sum_{\alpha=1}^3\sum_{w,z\in\{\mathsf u,\mathsf v\}}
 |b_{p_n}^\alpha(w,z)|^2\right]\le C.
\end{equation}
All constants are uniform on \eqref{eq:hierarchycone}.
\end{proposition}
\begin{proof}
For a parallel $k$, both $\operatorname{ad}k$ and
$I_\alpha\operatorname{ad}k$ preserve reference harmonic one-forms
and their orthogonal complement. Thus mixed harmonic/complementary
pairings vanish in \eqref{eq:vertices}. Lemma~\ref{lem:actualstates}
leaves an $O(r^2)$ error: the harmonic parts change by $O(r^2)$ and
both complementary parts are $O(r)$. Lemma~\ref{lem:lightdirection}
and spatial multiplication estimates allow replacement of each low
profile by its parallel approximation, with the same error.

At the reference plane, Pauli commutators give
$j_{k_3}=-2/\sqrt{6V_0}$ and $j_{k_1}=j_{k_2}=0$. The three flat
moment maps vanish identically on the Hermitian $X_3$ plane. Their
Hessians on each selected pair therefore vanish for every parallel
generator. This proves \eqref{eq:lowvertices}. The persistent central
and spectator generators commute pointwise with the selected
$\End E$ fluctuations, proving the exact zero statements.

Each functional in \eqref{eq:vertices} is the $L^2$ pairing of $p$
with a pointwise bilinear zero-form expression in $\mathsf u,\mathsf v$.
Lemma~\ref{lem:actualstates} and Sobolev multiplication bound its
$L^2$ norm uniformly. Parseval gives \eqref{eq:parsevalvertices},
including all complementary $\SO(32)$ blocks.
\end{proof}
